\newpage
\ifdefined\useChapters
\chapter{Como chegamos até aqui}
\else
\chapter{}
\fi

Desde que tudo isso começou, já senti tanta coisa, já vivi tanta coisa, mas nunca me canso de me surpreender. Estou o dia todo com uma voz dentro da minha cabeça, como se um pensamento fosse colocado lá dentro sem meu consentimento, e o mais estranho de tudo isso é que as duas pessoas que conheço que poderiam estar fazendo isso estão agora mortas. Eu mesma dei conta deles. Não tive muito sucesso recrutando aquela menina insuportável, talvez Pedro tenha tido mais sorte.

— Pedro, onde você está? — Fala Larissa, enquanto coloca a mão na cabeça, sentindo como se tivesse alguém pressionando.

— Então, Larissa, eu tenho uma notícia ruim pra te dar.

— Que droga, Pedro, você não conseguiu trazer o rapaz novo pro nosso lado, não é?

— O garçom chegou antes e o pior não é só isso, Jonas foi com ele também.

— Caramba, Pedro, dessa vez você se superou. Me encontra lá no lugar de sempre.

— Mas e você, conseguiu alguma coisa?

— Isso não vem ao caso, Pedro, aliás estamos perdendo tempo aqui, até mais.

Larissa, como filha adotiva de Crispim, cresceu sabendo que não podia confiar em ninguém e não podia esperar as situações piorarem pra tomar uma providência, porque quando menos você espera elas podem tomar proporções incontroláveis.

O fim da tarde no centro da cidade próximo ao coreto, desde que o corpo da mulher de vestido vermelho foi lá encontrado, é muito pouco movimentado e justamente por isso passou a ser o local e horário preferidos para um encontro sem serem observados.

— Eu sei que sempre te pergunto isso, mas os últimos fatos me mostraram que não posso confiar tanto em você, então me diz, você foi seguido por alguém? — Pergunta Larissa, levantando da mureta do coreto enquanto Pedro se aproximava.

— Você é sempre a senhora perfeitinha, não é mesmo? E sempre irritada, não é? — Fala Pedro, aproximando-se de Larissa com uma feição de derrotado. — Não, Larissa, ninguém me seguiu, eu dei algumas voltas nos quarteirões pra ter certeza que ninguém estava atrás de mim.

Pedro uma vez viu em um filme de detetives que se você ficar mudando de direção aleatoriamente por um tempo, quebra o planejamento de alguém que estiver te seguindo, fazendo com que a pessoa desista. Mas ele não tinha parado para pensar que não era só ele que assistia a esses filmes da sessão da tarde.

Ao longe, o garçom retornava caminhando junto a Bento e Jonas, combinando como fariam para acabar de uma vez com as chances de Crispim e os outros terem todo o poder que ele julgava ser somente seu por direito. Como ele mesmo havia matado a jovem no coreto, não havia por que ter receio de passar por aquela região, muito pelo contrário. Residindo em seu interior, havia um certo prazer em retornar aos locais de seus assassinatos, como se sentisse novamente toda a adrenalina e prazer que havia sentido no momento.

Enquanto caminhavam pelo centro, não pôde deixar de notar Pedro andando feito um doido de um lado para outro em direção ao coreto, logo percebendo que alguma coisa importante ele estava indo fazer. Pedro nunca foi muito esperto, era um rapaz com boas intenções e muito determinado, mas também, por sempre acreditar que podia resolver tudo com a força, não gastava muito tempo planejando suas ações.

Acima do coreto, enquanto Pedro e Larissa conversavam e o garçom e os outros se aproximavam, desciam Ítalo e o garoto, volitando bem devagar, tentando se aproximar sem serem notados.

— Como eu imaginei vocês iam estar se encontrando aqui hoje — Disse o garoto, enquanto atravessavam o teto do coreto, fazendo com que todos ficassem paralisados por um tempo.

— O que você está fazendo aqui? Eu tinha te matado junto com aquele imprestável do Crispim — fala Larissa, acendendo as chamas nas suas mãos.

A última conversa com Crispim tinha colocado no garoto um pouco mais de determinação, e sua última visita ao futuro fez com que tivesse muito mais crença em si mesmo. O que fez com que suas habilidades aumentassem muito, de maneira que, assim que chegou próximo de Larissa, segurou suas mãos em chamas e nada lhe aconteceu.

— Eu não estou aqui pra lutar com ninguém, até porque não deveríamos utilizar as nossas habilidades para machucar uns aos outros — fala o garoto, enquanto faz com que as chamas das mãos de Larissa desapareçam.

— Que droga, como você fez isso? — Pergunta Larissa, empurrando o garoto.

— Eu sempre me surpreendo como, mesmo você vendo tantas coisas, continua ainda duvidando que algo é possível.

— Nós precisamos nos unir, eu vi o futuro e não podemos permitir que ele aconteça. 

— Agora você virou vidente — fala Pedro, se aproximando dos dois com algumas pedras suspensas ao seu redor.

— Pedro, dá um tempo com essas pedras, vamos ouvir o que ele tem pra falar — Fala Larissa, se aproximando de Pedro.

— Eu sei que vocês ainda não confiam nele, mas em mim podem confiar. Estamos nessa há um tempo e eu ainda não dei nenhum motivo pra duvidarem. Tenho andado com ele ultimamente e, apesar de eu não ter visto esse futuro que ele diz, vi do que ele é capaz e acredito que ele esteja falando a verdade — fala Ítalo, chegando próximo de Pedro.

Por um tempo, o garoto fica quieto olhando para os quatro, mas não fala coisa alguma. Aquela tarde fria e escura de inverno fazia com que tudo ficasse um tanto mais sombrio.

— Eu imaginava que ia demorar um pouco mais pra esse confronto acontecer, mas agora já não podemos fugir — fala o garoto, colocando-se em posição de batalha e olhando para trás de Pedro e Larissa.

— O que está acontecendo, não estou vendo nada, com quem você está falando? — Fala Pedro, virando assustado para a mesma direção que o garoto olhava, mas sem entender nada.

— Preparem-se — fala o garoto, enquanto começa a reluzir como se estivesse envolto por uma fumaça brilhante. 

— O garçom e mais duas pessoas estão vindo, consigo sentir três pessoas, mas ainda não reconheço suas energias.

Do meio de alguns arbustos no fundo do coreto surgem voando um banco de madeira e algumas pedras grandes. Todos pulam para se defender, mas na verdade tudo não passava de uma distração. Do outro lado do coreto estava ele, com seus cabelos grisalhos e os bigodes queimados, o garçom, caminhando junto a dois velhos conhecidos do garoto, Bento e Jonas. 

Larissa e Pedro partem para cima do garçom, contando com uma batalha corpo a corpo, pois o garçom tinha a habilidade de fazer com que outros não tivessem habilidades especiais, mas dessa vez foi diferente. O garçom partiu para cima com as mãos em chamas, assim como Larissa fazia, lançando labaredas por todos os lados.

Bento e Jonas partem para cima do garoto, mas ele não revida, começando a fugir em volta do coreto. Ítalo achava tudo aquilo muito estranho: por que ele estava fugindo se podia muito bem acabar com tudo aquilo sem esforço algum?

— Ítalo, atraia todo mundo para cá, em frente ao coreto — transmite o garoto uma mensagem diretamente na mente, enquanto foge dos dois que tentam atingi-lo com pedras arrancadas da praça.

Ítalo passa então volitando, atraindo a atenção de todos. Ao mesmo tempo, o garoto envia a mesma mensagem para Larissa e Pedro. Todos se encaminham para a frente do coreto.

Quando todos estão em frente ao coreto, o garoto aparece no meio de todos e, com um grande clarão de luz, faz com que fiquem paralisados.

— Eu vou acabar com tudo isso agora, me desculpem. Não queria ter que fazer isso, mas vou levar vocês para muito tempo atrás, o mais longe que eu puder, para quando vocês não vão poder fazer mal pra ninguém.

Uma grande fumaça luminosa envolve todos em frente ao coreto e logo não havia ali mais ninguém. Todos tinham sido transportados para bem longe, mas diferente do que o garoto imaginava, não estavam em um momento diferente da história, talvez com alguns dinossauros como ele imaginava, mas sim no centro de Nova Iorque nos dias de hoje.

Transportar tantas pessoas era muito mais cansativo do que ele imaginava e logo cai desmaiado no centro da Times Square.

— O que é isso? Parece que não deu muito certo, não é, seu garoto idiota? Aliás, muito obrigado por ter me trazido para o lugar perfeito para o que eu quero fazer, onde mais ter tanta visibilidade quanto aqui no centro comercial do mundo.

— Ei, garoto, levanta, não deu muito certo levar as pessoas para onde você queria — fala Ítalo, balançando o garoto nos braços.

— Eu não estou entendendo, até parece que vocês estão facilitando as coisas pra mim — fala o garçom, enquanto corre para cima do garoto com labaredas nas mãos e um grande sorriso debochado no rosto.

Rapidamente, Ítalo pega o garoto no colo, como já havia feito antes, e dispara no meio dos prédios, desviando de algumas labaredas e latões de lixo arremessados. Bento

 e Jonas eram muito persistentes e os perseguiam por muitas ruas. O garoto era um pouco pesado e diminuía consideravelmente a fuga dos dois e, enquanto volitavam pelo centro de Nova Iorque, Ítalo e o garoto eram feridos, o que fazia com que cada quarteirão ficasse ainda mais difícil despistar.

Alguns momentos críticos fazem com que alianças sejam criadas para fazer frente a um inimigo comum e, apesar de Larissa ser muito reticente a se aliar ao garoto porque ele no princípio lhe trouxe muitos problemas e mesmo agora tentara prendê-la no passado, também sabia que alguém como o garçom com mais poder do que deveria ter poderia causar um dano talvez até irreversível.

Enquanto garçom e os dois os perseguiam, não puderam notar que atrás deles estavam Larissa e Pedro, que também com várias labaredas defendiam quem eles nem imaginavam.

— Eu não sei, Ítalo, como vamos fazer para reverter o que eu vi — fala o garoto, começando a despertar.

— Você está acordado? Não gaste tanta energia, você se esforçou muito agora pouco.

— Apesar de ter voltado no tempo e de ter tentado mudar tantas coisas, tudo parece convergir sempre pro mesmo futuro. Então foi assim que todos vieram para Nova Iorque, foi assim que eu tinha visto tudo acontecer aqui, parece ser tudo inevitável.

— É, eu realmente não sei mesmo o que podemos fazer, mas também não acho que por isso temos que desistir. Pode até ser que não consigamos evitar tudo o que você viu, mas também pelo que você me conta não era o fim, então podemos evitar que dali para frente as coisas caminhem para um futuro horrível — fala Ítalo, com a sua voz compassiva de sempre.

— É verdade, Ítalo, eu não tinha pensado nisso. Vamos então fazer com que o futuro mesmo seja algo bom.

